%%%%%%%%%%%%%%%%%%%%%%%%%%%%%%%%%%%%%%%%%%%%%%%%%%%%%%%%%%%%%%%%%%%%%%%%%%%%%%%%
% introducao.tex
%
% Capítulo de introdução da monografia.
%%%%%%%%%%%%%%%%%%%%%%%%%%%%%%%%%%%%%%%%%%%%%%%%%%%%%%%%%%%%%%%%%%%%%%%%%%%%%%%%

\chapter{Introdução}
\label{sec:intro}

A transformação digital no setor público tem impulsionado a revisão de
processos administrativos e educacionais, especialmente em instituições cuja
atuação depende de controle operacional, rastreabilidade das informações e
padronização de rotinas. No contexto da formação profissional, sistemas
integrados tornam-se relevantes por apoiar atividades como planejamento de
cursos, gestão de inscrições, acompanhamento de participação, registro de
frequência, certificação e demais etapas administrativas que compõem o ciclo
pedagógico.

Nesse cenário, escolas da magistratura desempenham papel estratégico na
qualificação de magistrados, servidores e demais públicos externos ou vinculados ao Poder
Judiciário, o que demanda instrumentos tecnológicos compatíveis com a
complexidade de seus fluxos institucionais. A Escola da Magistratura do Espírito
Santo (EMES) insere-se nesse contexto ao reunir atividades acadêmicas,
administrativas e comunicacionais que precisam operar de forma articulada,
segura e eficiente.

Com o objetivo de apoiar essas demandas, foi desenvolvido e lançado em novembro
de 2025 o sistema Educa EMES, uma solução integrada de gestão educacional
construída em contexto real de aplicação institucional e atualmente utilizada
pela EMES. O sistema passou a concentrar diferentes módulos relacionados à
gestão acadêmica e administrativa, oferecendo suporte a processos como
inscrição em cursos, controle de frequência, emissão de certificados,
mensageria, documentação e operações internas da instituição.

Embora o Educa EMES represente um avanço na modernização e integração dos
processos da EMES, o fluxo pedagógico ainda apresenta etapas com intervenção
manual, ausência de padronização em determinadas rotinas e limitações quanto à
automação integral do ciclo acadêmico. Diante disso, este trabalho não trata do
desenvolvimento do sistema Educa EMES como um todo, mas da análise e modelagem
de uma solução de automação do fluxo pedagógico dentro de uma plataforma já
implantada. Nessa perspectiva, o Educa EMES é tratado simultaneamente como
ambiente base da solução proposta e como estudo de caso aplicado à gestão
educacional em instituição pública.

\section{Descrição do problema}
\label{sec:intro:prob}

Antes da adoção de uma solução integrada, a EMES executava parte relevante de
seus processos pedagógicos por meio de procedimentos manuais. Entre essas
rotinas estavam registros físicos de frequência, controles operacionais
dispersos e atividades dependentes de intervenção direta dos servidores, o que
reduzia a eficiência das ações pedagógicas e dificultava a consolidação das
informações acadêmicas.

Em 2023, a instituição passou a utilizar o sistema Emeron Web como alternativa
de apoio ao fluxo pedagógico. Contudo, a solução apresentava limitações
relacionadas à usabilidade, à adequação arquitetural e ao escopo funcional,
atendendo apenas parte das necessidades da escola. Como consequência, diversas
atividades continuaram sendo realizadas de maneira manual, como comunicação com
alunos, envio de mensagens, divulgação de cursos e aplicação de formulários.

A implementação do Educa EMES representou um avanço significativo ao integrar
processos antes fragmentados e fornecer uma base tecnológica mais aderente ao
contexto institucional. Ainda assim, verificou-se que a existência de um sistema
em produção não eliminou completamente os pontos de intervenção manual no fluxo
pedagógico. Persistem lacunas relacionadas à automação de etapas, à
padronização de procedimentos e à redução da dependência de operações
executadas manualmente ao longo do ciclo acadêmico.

Desse modo, o problema abordado neste trabalho não consiste na inexistência de
um sistema para a EMES, mas na necessidade de ampliar a automação do fluxo
pedagógico dentro de uma plataforma já implantada. Essa necessidade envolve
identificar gargalos e limitações do fluxo atual, propor melhorias baseadas em
automação e modelar mecanismos capazes de aumentar a eficiência operacional, a
rastreabilidade e a consistência dos procedimentos acadêmicos.

\section{Formulação do problema}
\label{sec:intro:form_prob}

Como modelar e desenvolver uma solução capaz de ampliar a automação do fluxo
pedagógico no sistema Educa EMES, reduzindo intervenções manuais, preservando a
rastreabilidade das etapas e ampliando a padronização dos processos acadêmicos
da EMES?

\section{Proposta de solução}
\label{sec:intro:prop_sol}

Como resposta ao problema identificado, este trabalho propõe a evolução do
sistema Educa EMES por meio da análise, da modelagem e da definição de
mecanismos voltados à automação do fluxo pedagógico. A proposta considera o
sistema já existente como base tecnológica da solução, concentrando-se
principalmente nos processos de inscrição, acompanhamento de frequência,
aprovação, emissão de certificados e geração de relatórios.

A proposta adotada corresponde a uma pesquisa aplicada, uma vez que a solução é
concebida a partir de uma demanda institucional concreta e analisada em um
ambiente real de uso. A contribuição do trabalho está na identificação de
lacunas e gargalos do fluxo atual, na definição de requisitos, na modelagem de
uma solução aderente às necessidades da EMES e no desenvolvimento de mecanismos
voltados à automação do fluxo pedagógico. A avaliação qualitativa dos impactos
da solução considera sua aderência ao contexto institucional, a redução de
intervenções manuais e a preservação da rastreabilidade das etapas.

Assim, o Educa EMES não é tratado como resultado final e acabado, mas como uma
plataforma integrada em processo de melhoria contínua. A solução proposta busca
consolidar o fluxo pedagógico dentro do sistema, ampliar o nível de automação
das rotinas acadêmicas e reduzir a dependência de procedimentos manuais ainda
existentes na operação institucional, preservando a possibilidade de
intervenção humana quando necessária.

\section{Objetivos}
\label{sec:intro:obj}

\subsection{Objetivo geral}
\label{sec:intro:obj:ger}

Analisar, modelar e desenvolver uma solução de automação do fluxo pedagógico no
sistema Educa EMES.

\subsection{Objetivos específicos}
\label{sec:intro:obj:esp}

\begin{itemize}
\item Analisar o fluxo pedagógico atual da EMES no contexto de utilização do
sistema Educa EMES.

\item Identificar os pontos de intervenção manual e as limitações de automação
existentes nas etapas do processo acadêmico.

\item Identificar gargalos e oportunidades de melhoria no fluxo pedagógico da
EMES.

\item Analisar a estrutura tecnológica e o modelo de dados do sistema atual no
recorte relacionado ao fluxo pedagógico.

\item Definir os requisitos necessários para a automação do fluxo pedagógico.


\item Modelar a solução proposta, contemplando etapas, estados, atores,
entidades e componentes arquiteturais.

\item Desenvolver mecanismos voltados à automação de etapas do fluxo pedagógico
no sistema Educa EMES.

\item Avaliar qualitativamente a solução proposta quanto à eficiência,
padronização, rastreabilidade e aderência ao contexto institucional.

\end{itemize}

\section{Justificativa}
\label{sec:intro:jus}

A proposta justifica-se pela relevância da automação de processos
educacionais em instituições públicas que dependem de fluxos administrativos e
acadêmicos confiáveis, integrados e rastreáveis. Nesses contextos, a
persistência de atividades manuais tende a gerar retrabalho, ampliar a
possibilidade de inconsistências e dificultar o acompanhamento das
ações realizadas ao longo do fluxo pedagógico.

No caso da EMES, a adoção do Educa EMES representou uma melhoria expressiva na
digitalização dos processos institucionais, mas também evidenciou que a
existência de um sistema integrado não encerra, por si só, as necessidades de
evolução tecnológica da organização. A permanência de etapas dependentes de
intervenção manual demonstra a importância de investigar como a automação pode
ser ampliada de maneira estruturada, preservando aderência ao contexto
institucional e promovendo ganhos de eficiência e padronização.

Sob a perspectiva acadêmica, a proposta contribui para a área de Sistemas de
Informação ao abordar uma experiência real de evolução de software em produção,
com ênfase na análise e modelagem de mecanismos de automação aplicados à gestão
educacional. Além disso, o recorte adotado nos processos de inscrição,
frequência, aprovação, certificação e geração de relatórios concentra-se em
etapas centrais do fluxo pedagógico, o que reforça a pertinência do estudo de
caso e sua potencial utilidade para contextos institucionais semelhantes.
